\subsection{The Core Language} \label{subsec:the-core-language}
\par While the user works with the surface language, the core language (see figure \ref{fig:the-core-language}). 
The core language is depth-indexed, meaning that the depth of the core AST is part of its type. 
The depth-index indicates the nested depth of the list it will generate. 
A depth of 0 indicates that the result will be a singleton not a list. 
We use this feature to ensure that all expressions are well formed before runtime. 
Furthermore, we can put constraints on what depths are allowed for each constructor. 
For example, the $\concat$ operator is for concatenation which requires 2 lists. 
So we can ensure that its two arguments have the same depth, and that it is non zero.

\begin{figure}
\begin{center}
    \begin{align*}
       \core &\to\core '|_{\lab}'\core   \\
            &\OR\core*_{\lab}                              \\
            &\OR\core \cons\core                     \\
            &\OR\core \concat\core                    \\
            &\OR (\lab \Longleftarrow b) \asgnarrow\core   \\
            &\OR \lab                                               \\        
            &\OR <x>                                                \\
            &\OR \nil                                               \\
    \end{align*}
\end{center}
    \caption{The Core Language as a Context Free Grammar}
    \label{fig:the-core-language}
\end{figure}